\documentclass[conference]{IEEEtran}
\usepackage{cite}
\usepackage{amsmath,amssymb,amsfonts}
\usepackage{algorithmic}
\usepackage{graphicx}
\usepackage{textcomp}
\usepackage{xcolor}
\usepackage{booktabs}
\usepackage{subcaption}
\usepackage{hyperref}

\def\BibTeX{{\rm B\kern-.05em{\sc i\kern-.025em b}\kern-.08em
    T\kern-.1667em\lower.7ex\hbox{E}\kern-.125emX}}

\begin{document}

\title{Beyond Static Policies: Exploring Dynamic Policy Selection for Single-Thread Performance Optimization}

\author{\IEEEauthorblockN{Yanxin Zhang, Ian McDougall, Junnan Li, Shayne Wadle, Vikas Singh, Karthikeyan Sankaralingam}
\IEEEauthorblockA{\textit{Department of Computer Sciences} \\
\textit{University of Wisconsin--Madison}\\
Madison, WI, USA \\
\{yanxin, imcdougall, junnan, swadle, vsingh, karu\}@cs.wisc.edu}}

\maketitle

\begin{abstract}
For over a decade, processor design has focused on implementing sophisticated policies for various components of the out-of-order pipeline, including cache replacement, prefetching, and branch prediction. The prevailing design philosophy has been to build processors with a single, static selection of policies across these different mechanisms. This paper investigates a fundamental question: do different workloads, or even different execution phases within the same workload, benefit from different policy combinations? We present a comprehensive analysis exploring whether a hypothetical processor capable of dynamically selecting from multiple policies could significantly outperform traditional static-policy processors. Using ChampSim-based simulation across 49 benchmarks segmented into 4,900 instruction chunks, we evaluate performance across multiple policy combinations for cache replacement, prefetching, and branch prediction. Our findings reveal that significant performance headroom exists---even the best static policy achieves optimal performance for only 38.67\% of execution phases and incurs a mean IPC loss of 1.38\% compared to an oracle. Furthermore, we demonstrate that a processor capable of dynamically switching between just two carefully chosen policies can achieve a 20.2$\times$ reduction in IPC loss (from 1.38\% to 0.068\%) and match oracle performance 70.43\% of the time. These results suggest that dynamic policy selection represents a promising avenue for unlocking single-thread performance improvements that have become increasingly difficult to achieve.
\end{abstract}

\begin{IEEEkeywords}
processor architecture, dynamic policy selection, cache replacement, prefetching, branch prediction, performance optimization
\end{IEEEkeywords}

\section{Introduction}

The past decade of processor microarchitecture has been characterized by the development and refinement of increasingly sophisticated policies embedded throughout the out-of-order execution pipeline. These policies govern critical architectural decisions including cache replacement strategies, hardware prefetching mechanisms, and branch prediction algorithms. Each of these components plays a crucial role in determining overall performance, yet the prevailing design philosophy has been to configure processors with a single, static combination of policies that remains fixed throughout program execution.

This design approach raises a fundamental question that has received limited attention in the architecture community: how much performance headroom exists if processors can dynamically select among multiple policies across execution phases?

This work provides a comprehensive answer to these questions through extensive simulation-based analysis. We explore whether dynamic policy selection represents a fruitful research direction---a new knob by which we can unlock single-thread performance improvements that have proven extremely difficult to achieve through conventional means.

Our key contributions include:

\begin{itemize}
\item A unified simulation framework integrating multiple state-of-the-art policies for L1 data cache prefetching, L1 instruction cache prefetching, and L2 cache replacement
\item Comprehensive performance analysis across 49 diverse benchmarks, each segmented into 100 execution phases (timesteps) of 200K instructions, totaling 4,900 distinct execution scenarios
\item Quantification of the performance gap between static policy selection and an oracle with perfect knowledge: the best static policy achieves only 38.67\% oracle match rate; IPC loss can exceed 10\% for many execution phases, while 26\% of execution phases show potential for $>1\%$ improvement through dynamic selection \textbf{alone}.
\item Demonstration that a two-policy dynamic selection mechanism can achieve a 20.2$\times$ reduction in mean IPC loss, while increasing oracle match rate from 38.67\% to 70.43\% of timesteps.
\item Analysis of policy selection patterns revealing heavy concentration on two dominant policies, with 63.49\% of timesteps already within 0.5\% of optimal using the best static choice
\end{itemize}

Our findings demonstrate that significant performance headroom exists for dynamic policy selection mechanisms. Even the best static policy is optimal for only 38.67\% of execution phases, incurring measurable IPC loss across diverse workloads that varies substantially across execution phases. A simple two-policy dynamic approach captures most of this headroom, achieving a median loss of exactly 0\% (meaning over half of all execution phases achieve perfect oracle-level performance) and a mean loss of just 0.068\%, while keeping 96.29\% of timesteps within 0.5\% of optimal. These results suggest that dynamic policy selection deserves serious
consideration as a practical mechanism for future performance
improvements in single-threaded execution. Designing an efficient
runtime mechanism to select among policies is an important direction
for future work and is beyond the scope of this paper.

\section{Methodology}

\subsection{Simulation Infrastructure}

Our analysis is built upon ChampSim\footnote{\url{https://github.com/ChampSim/ChampSim}}, a widely-used trace-driven microarchitecture simulator that has been extensively validated against real hardware~\cite{gober2022championship}. ChampSim provides a flexible framework for implementing and evaluating various architectural policies while maintaining simulation speed sufficient for large-scale studies.

We extended ChampSim to implement multiple state-of-the-art policies within a unified codebase. This unified implementation is critical for ensuring fair comparison, as it eliminates variability that could arise from different simulator configurations, warmup periods, or measurement methodologies. All policies were implemented with careful attention to published specifications and validated against reported results where possible.

\subsection{Policy Implementations}

Our study evaluates policies across three key microarchitectural components. 
For L1 data cache prefetching, we implement Berti~\cite{navarro2022berti} and Gaze~\cite{chen2025gaze}. 
For L1 instruction cache prefetching, we evaluate Entangling~\cite{ros2021cost} and BARCA~\cite{gratz2020barca}. 
For L2 cache replacement, we implement Mockingjay~\cite{shah2022effective} and PACIPV~\cite{mostofi2025light}.
\subsection{Benchmark Suite}

We selected 49 diverse benchmarks representing a wide range of application domains and execution characteristics. These benchmarks were chosen to provide comprehensive coverage of different computational patterns, memory behaviors, and control flow characteristics.

\subsection{Temporal Segmentation}

A key aspect of our methodology is the temporal segmentation of each benchmark into discrete execution phases. Each benchmark trace was divided into chunks of 200,000 instructions, which we refer to as \textit{timesteps}. This granularity was chosen based on several considerations:

\begin{itemize}
\item It is fine enough to capture phase behavior within applications
\item It is coarse enough to amortize any overhead of policy switching
\item It provides sufficient instructions for stable IPC measurements
\end{itemize}

With 49 benchmarks and 100 timesteps per benchmark, our complete evaluation spans 4,900 distinct execution phases, providing a comprehensive dataset for analysis.

\subsection{Performance Metrics}

For each combination of benchmark, timestep, and policy configuration, we measure Instructions Per Cycle (IPC) as the primary performance metric. 
All IPC values are collected across the entire execution timeline for 
every application and policy configuration under study.

From these raw IPC measurements, we derive several analysis metrics:

\textbf{IPC Loss:} For each timestep, we compute the IPC loss of a static policy relative to the oracle (the best-performing policy at that timestep):
\begin{equation}
\text{Loss} = \frac{\text{IPC}_{\text{oracle}} - \text{IPC}_{\text{policy}}}{\text{IPC}_{\text{oracle}}} \times 100\%
\end{equation}

\textbf{Oracle Performance:} The oracle represents an idealized processor with perfect knowledge of which policy will perform best at each timestep. While not practically realizable, it provides an upper bound on the performance achievable through dynamic policy selection.

\textbf{Baseline Performance:} We define a baseline processor configuration and report its performance under a fixed static policy selection; at minimum, this baseline is specified by issue width, ROB size, and L1/L2 cache sizes.


\section{Results}

\subsection{Single Policy Analysis: Per-Application Variability}

We first analyze how policy effectiveness varies across execution
phases within individual applications. 
Figure~\ref{fig:loss_dist_single} shows the IPC loss distribution of a
representative static policy (Berti/Entangling/Mockingjay) across all
benchmarks.

The results reveal substantial temporal variability: even with the best-by-mean static policy (berti/barca/PACIPV), performance varies widely across execution phases, indicating strong potential benefit from dynamic policy selection. Across all 4,900 phases, IPC loss exceeds 10\% in 210 phases (4.29\%), exceeds 5\% in 441 phases (9.00\%), and exceeds 1\% in 1,280 phases (26.12\%).

\begin{figure}[tb]
\centering
\includegraphics[width=0.48\textwidth]{loss_distribution_Berti_Entangling_Mockingjay.pdf}
\caption{Box plot showing the distribution of IPC loss across timesteps for each benchmark when using a single static policy (Berti/Entangling/Mockingjay). Each box represents the distribution across 100 timesteps for one benchmark. Wide distributions indicate high temporal variability, suggesting that different execution phases would benefit from different policies.}
\label{fig:loss_dist_single}
\end{figure}

\subsection{Global Static Policy Performance and Policy Comparison}

To quantify overall performance headroom, we aggregate all 4,900 timesteps
(49 benchmarks $\times$ 100 timesteps) and treat each timestep as an
independent data point. Figure~\ref{fig:static_vs_oracle_single} shows the
IPC loss distribution of a representative static policy relative to the
oracle.

\begin{figure}[tb]
\centering
\includegraphics[width=0.48\textwidth]{static_vs_oracle_dist_Berti_Entangling_Mockingjay.pdf}
\caption{Global distribution of IPC loss for a representative static policy compared to the oracle across all 4,900 timesteps. Each data point corresponds to one 200K-instruction execution phase.}
\label{fig:static_vs_oracle_single}
\end{figure}

Even the best-performing static configuration (Berti/BARCA/PACIPV)
achieves only a 38.67\% oracle match rate and incurs a mean IPC loss
of 1.38\%. Moreover, 26.12\% of timesteps show potential improvement
greater than 1\%, indicating substantial performance headroom for
dynamic policy selection.

To determine whether this headroom is specific to one configuration,
Figure~\ref{fig:all_policies} compares IPC loss distributions across
all evaluated static policy combinations.

\begin{figure*}[tb]
\centering
\includegraphics[width=0.95\textwidth]{static_vs_oracle_dist_summary.pdf}
\caption{Comparison of IPC loss distributions across all evaluated static policy combinations. Each distribution is measured relative to the oracle across all 4,900 timesteps.}
\label{fig:all_policies}
\end{figure*}

No single static configuration dominates across all timesteps, and
policy effectiveness varies substantially across execution phases,
highlighting the structural limitations of static policy selection.

\subsection{Policy Optimality and Two-Policy Dynamic Selection}

To examine whether oracle optimality is concentrated on a small number
of policies, we analyze how frequently each static configuration
achieves the highest IPC at each timestep.
Figure~\ref{fig:policy_freq} shows the optimality frequency across
all 4,900 timesteps.

\begin{figure}[tb]
\centering
\includegraphics[width=0.48\textwidth]{policy_optimality_frequency.pdf}
\caption{Frequency with which each policy configuration is optimal across all 4,900 timesteps.}
\label{fig:policy_freq}
\end{figure}

Optimality is heavily concentrated on two dominant policy
combinations. The most frequently selected winner is
Gaze/Entangling/Mockingjay at 34.94\% of timesteps, followed by
Berti/Entangling/Mockingjay at 34.35\%; together, these two
configurations cover 69.29\% of timesteps.

Motivated by this observation, we evaluate a processor capable of
dynamically switching between two policies. Figure~\ref{fig:two_policy}
shows the resulting IPC loss distribution relative to the oracle.

\begin{figure}[tb]
\centering
\includegraphics[width=0.48\textwidth]{policy_vs_baseline_gaze_barca_PACIPV.pdf}
\caption{IPC loss distribution for a processor capable of dynamically selecting between two dominant policies.}
\label{fig:two_policy}
\end{figure}

Compared to the best static configuration (mean IPC loss 1.38\%),
the two-policy mechanism reduces mean loss to 0.068\%
(20.2$\times$ improvement) and increases the oracle match rate
from 38.67\% to 70.43\% of timesteps.

These results show that supporting only two selectable policies
captures most of the oracle's performance advantage while avoiding
the complexity of full dynamic selection.

\subsection{Summary of Key Findings}

Taken together, our results show that static policy selection leaves 
measurable performance headroom across diverse workloads and execution 
phases. This headroom is driven by substantial temporal variability 
within applications and a clear phase-dependent ranking among policy 
configurations. Importantly, optimality is heavily concentrated on a 
small number of policies, enabling a practical two-policy dynamic 
mechanism to capture most of the oracle’s performance advantage while 
avoiding the complexity of full dynamic selection. Although sensitivity 
varies across workloads, the overall pattern consistently indicates 
that dynamic policy selection represents a viable and structurally 
justified direction for improving single-thread performance.

\section{Related Work}

Prior work has explored adaptive mechanisms across several
microarchitectural components. Adaptive cache replacement
policies such as DRRIP adjust replacement decisions based on
workload behavior, but operate within a single policy framework.

Adaptive prefetching research has shown that prefetcher
effectiveness varies across applications and execution phases,
with mechanisms such as throttling and confidence-based
control.

Branch prediction has also explored hybrid predictors that
combine multiple prediction strategies, though typically within
the branch prediction unit itself.

Prior work on phase detection and reconfigurable architectures
has further investigated runtime adaptation of architectural
behavior. In contrast, our work quantifies the performance
opportunity of dynamically selecting among multiple
microarchitectural policies across execution phases.

\section{Discussion}

While our results quantify substantial performance headroom, practical
realization of dynamic policy selection requires addressing switching
overhead and phase detection mechanisms.

Our analysis shows that optimality is strongly concentrated on two
dominant policies, suggesting that a simple two-way selection mechanism
can capture most of the oracle’s benefit without supporting full
dynamic selection. We do not claim this to be universal; depending on
the policy space considered, more than two policies may be required.

These observations indicate that dynamic policy selection represents a
practical and architecturally tractable opportunity for improving
single-thread performance.
\section{Conclusion}

We quantify the performance headroom left by static
microarchitectural policy selection across 49 benchmarks and
4,900 execution phases. Our results show that optimal policy
choice varies across phases and that no single static
configuration dominates.

Because optimality is concentrated on two dominant policies,
a simple two-policy dynamic mechanism captures most of the
oracle advantage, reducing mean IPC loss from 1.38\% to 0.068\%
and achieving oracle-level performance in over 70\% of phases.

These findings highlight dynamic policy selection as a practical
direction for improving single-thread performance.

\section*{Acknowledgments}

The authors would like to thank the developers of ChampSim for providing an excellent simulation infrastructure that made this research possible.

\bibliographystyle{IEEEtran}
\bibliography{refs}

\end{document}
